\documentclass[aspectratio=43,11pt]{beamer}
\usetheme{default}
\usecolortheme{default}
\setbeamertemplate{navigation symbols}{}
\setbeamertemplate{footline}[frame number]
\setbeamertemplate{blocks}[rounded][shadow=false]
\setbeamersize{text margin left=12mm,text margin right=12mm}
\setlength{\parskip}{4pt}
\newcommand{\visual}[3]{\begin{frame}{#1}\centering\includegraphics[width=.94\textwidth,height=.71\textheight,keepaspectratio]{#2}\par\vspace{1mm}{\small #3}\end{frame}}
\usepackage[T1]{fontenc}
\usepackage{lmodern}
\usepackage{amsmath,amssymb,booktabs,listings,xcolor,graphicx}
\definecolor{aggiemaroon}{RGB}{80,0,0}
\setbeamercolor{structure}{fg=aggiemaroon}
\setbeamercolor{frametitle}{fg=aggiemaroon}
\setbeamercolor{block title}{bg=black!8,fg=black}
\setbeamercolor{block body}{bg=black!3,fg=black}
\definecolor{codegray}{gray}{0.24}
\lstdefinestyle{R}{language=R,basicstyle=\ttfamily\small,keywordstyle=\color{black},commentstyle=\color{codegray},showstringspaces=false,breaklines=true,columns=fullflexible,keepspaces=true}
\lstset{style=R}
\newcommand{\stage}[2]{\begin{block}{#1}\small #2\end{block}}
\title{Decision and Risk Analysis}
\subtitle{Lecture 11: \textbf{Normal Hierarchical Regression\\with Predictors}}
\author{}
\date{\color{aggiemaroon}{DAEN 427 / ISEN 427 / ISEN 627 - Fall 2026\\
      \textbf{Alexander Abuabara}}}
\titlegraphic{\vfill \includegraphics[height=1.2cm]{TAMU_logo.png}}

\begin{document}

\begin{frame}\titlepage
\centering \footnotesize \url{https://www.bayesrulesbook.com/chapter-17}
\end{frame}

\begin{frame}{}
\begin{enumerate}
\item \textbf{Evaluate:} critique a pooled age--running-time analysis.
\item \textbf{Apply:} explore groups and fit a varying-intercept model.
\item \textbf{Understand:} interpret hierarchy, shrinkage, priors, and slopes.
\item \textbf{Evaluate:} check diagnostics and compare models.
\item \textbf{Create:} design predictions and transfer the model to another setting.
\item \textbf{Evaluate:} defend model choice and limitations.
\end{enumerate}
\end{frame}

\begin{frame}{Does age really have so little effect?}
\stage{Case}{Cherry Blossom 10-mile race: repeated net times for \textbf{36 runners} in their 50s and 60s; 252 complete races reported in the chapter.}
\[Y_{ij}=\text{net minutes for race }i\text{ of runner }j,\quad X_{ij}=\text{age}.\]
\begin{block}{Complete pooling}\centering $Y_{ij}\sim N(\beta_0+\beta_1X_{ij},\sigma^2)$\end{block}
Book:\\
pooled age slope $\approx0.268$ min/year;
80\% CrIn $[-0.302,0.842]$.\par
Does a weak pooled trend imply runners do not slow down?\\
What assumption is hidden?
\end{frame}

\begin{frame}{MCQ 1 --- Diagnose complete pooling}
We draw one regression line through every recorded race. What important detail have we left out?
\begin{enumerate}
\item[A.] Runners cannot have different ages.
\item[B.] Repeated outcomes on a runner may be correlated.
\item[C.] Normal models cannot use numerical predictors.
\item[D.] Bayesian models cannot estimate slopes.
\end{enumerate}
\textcolor{gray}{\it Answers are at the end.}
\end{frame}

\begin{frame}[fragile]{Coding warm-up (inspect before modeling)}
\begin{lstlisting}
library(bayesrules)
library(dplyr)
library(ggplot2)
data(cherry_blossom_sample)
running <- cherry_blossom_sample |>
  select(runner, age, net) |> na.omit() |>
  mutate(runner = factor(runner))

nrow(running); n_distinct(running$runner)
running |> count(runner) |> summary()
ggplot(running, aes(age, net)) + geom_point() +
  facet_wrap(~ runner)
\end{lstlisting}
\textbf{Discuss:} What do you notice about the runners?\\
Look for differences in typical times and changes with age.
\end{frame}

\visual{A pooled line can be misleading}{ch17_pool.png}{Illustrative data: one line cannot describe each runner\'s history.}

\begin{frame}{Give each runner a starting point}
\begin{align*}
Y_{ij}\mid\beta_{0j},\beta_1,\sigma_y &\sim N(\beta_{0j}+\beta_1 X_{ij},\sigma_y^2),\\
\beta_{0j}\mid\beta_0,\sigma_0 &\sim N(\beta_0,\sigma_0^2).
\end{align*}
\stage{}{Each runner gets their own typical time, but we learn one common age trend. The runners help us learn about one another.}
$\sigma_y$: race-to-race scatter within a runner.\\
$\sigma_0$: between-runner scatter in intercepts.\par
\textbf{Sketch:}\\Draw five parallel runner regression lines and the global line.
\end{frame}

\visual{What varying intercepts look like}{ch17_intercepts.png}{Each line is one runner. Different baselines, but parallel slopes.}

\begin{frame}[fragile]{Hands-on: build complete and varying-intercept fits}
\begin{lstlisting}
library(rstanarm)
# Short teaching runs; increase iter for final analysis.
m_pool <- stan_glm(net ~ age, data = running,
  family = gaussian(), chains = 2, iter = 1000,
  seed = 84735, refresh = 0)

m_int <- stan_glmer(net ~ age + (1 | runner),
  data = running, family = gaussian(),
  prior_intercept = normal(100, 10),
  prior = normal(2.5, 1),
  prior_aux = exponential(1, autoscale = TRUE),
  prior_covariance = decov(reg=1, conc=1,
                           shape=1, scale=1),
  chains = 2, iter = 1000, seed = 84735,
  refresh = 0)
\end{lstlisting}
\textbf{Task:} Explain the \texttt{(1 | runner)} term; label the prior on each global parameter.
\end{frame}

\begin{frame}{MCQ 2 --- What changes from runner to runner?}
In \texttt{net \~{} age + (1 | runner)},\\
which feature changes between runners?
\begin{enumerate}
\item[A.] Intercept only
\item[B.] Age slope only
\item[C.] Both intercept and age slope
\item[D.] Neither intercept nor slope
\end{enumerate}
\end{frame}

\begin{frame}[fragile]{Hands-on: inspect and interpret the output}
\begin{lstlisting}
library(broom.mixed)
tidy(m_int, effects = "fixed", conf.int = TRUE,
     conf.level = 0.80)
tidy(m_int, effects = "ran_pars")
ranef(m_int)$runner |> head()
\end{lstlisting}
\stage{Book benchmark (longer fit)}{Global age slope: about $1.30$ min/year, with 80\% credible interval $[1.02,1.58]$. Between-runner SD $\approx13.3$ min; within-runner SD $\approx5.25$ min. Short class runs will differ.}
\textbf{Question:} Translate the age interval into a sentence about an individual runner's expected change per year, under the model.
\end{frame}

\begin{frame}{The three choices, in plain language}
\textbf{Complete pooling:} Treat all races as if they came from the same runner.
\par \vspace{5mm}
\textbf{No pooling:} Analyze every runner separately; learn nothing across runners.
\par \vspace{5mm}
\textbf{Partial pooling:} Let runners differ, while sharing information about the broader group.
\par \vspace{6mm}
\textbf{Ask:} Which approach would you trust for someone with only two races?
\end{frame}

\visual{Why share information?}{ch17_shrinkage.png}{Illustration: estimates based on fewer races usually move more toward the group.}

\begin{frame}{What do our priors say before seeing the data?}
\begin{align*}
\beta_{0c} &\sim N(100,10^2), & \beta_1&\sim N(2.5,1^2),\\
\sigma_y &\sim\mathrm{Exp}(0.072), & \sigma_0&\sim\mathrm{Exp}(1).
\end{align*}
The intercept prior applies to the \emph{centered} intercept, evaluated at mean observed age, in the book's \texttt{rstanarm} parametrization, not the time at age zero.\par
\par \vspace{4mm}
Is a mean age slope of 2.5 min/year plausible before examining the data? Why simulate from the prior?
\end{frame}

\begin{frame}[fragile]{Do these priors generate sensible times?}
\begin{lstlisting}
m_prior <- update(m_int, prior_PD = TRUE,
                  chains = 2, iter = 1000,
                  seed = 84735, refresh = 0)
pp_check(m_prior, plotfun = "dens_overlay", nsamples = 30)
pp_check(m_int, plotfun = "dens_overlay", nsamples = 30)
\end{lstlisting}
\textbf{Identify:} Any impossible negative times? Extremely slow 10-mile races? How should a prior predictive check affect prior choice?\par
\end{frame}

\begin{frame}{What if runners age differently?}
\[\mu_{ij}=\beta_{0j}+\beta_{1j}X_{ij},\qquad
\begin{pmatrix}\beta_{0j}\\\beta_{1j}\end{pmatrix}\sim N_2\!\left(\begin{pmatrix}\beta_0\\\beta_1\end{pmatrix},\Sigma\right)\]
\[\Sigma=\begin{pmatrix}\sigma_0^2&\rho\sigma_0\sigma_1\\\rho\sigma_0\sigma_1&\sigma_1^2\end{pmatrix}.\]
Runners can differ in both baseline time and age-related change. $\rho$ describes their intercept--slope association.\par
\textbf{Chapter example:} Runner 20 appears to slow more rapidly than runner 29 in the raw data. Is that sufficient evidence to fit separate slopes?
\end{frame}

\visual{What varying slopes look like}{ch17_slopes.png}{Here, some runners change more rapidly with age than others.}

\begin{frame}{MCQ 3 --- Interpret correlation}
Suppose the intercept--slope correlation is negative, where a higher intercept means a slower baseline time and a higher slope means greater slowing with age. Then:
\begin{enumerate}
\item[A.] Slower-baseline runners tend to slow down faster.
\item[B.] Faster-baseline runners tend to slow down faster.
\item[C.] Every runner has the same slope.
\item[D.] Race-to-race residual variation must be zero.
\end{enumerate}
\end{frame}

\begin{frame}{Try allowing different slopes}
\textbf{Goal:} Decide whether varying age slopes materially improve the analysis.\par
\vspace{3mm}
\begin{block}{Models compared}
\texttt{net \~{} age}\par
\texttt{net \~{} age + (1 | runner)}\par
\texttt{net \~{} age + (age | runner)}
\end{block}
The final formula estimates correlated runner-specific intercepts and slopes. Greater flexibility can increase computation and estimation difficulty.
\end{frame}

\begin{frame}[fragile]{Coding lab: varying slopes and diagnostics}
\begin{lstlisting}
m_slope <- stan_glmer(net ~ age + (age | runner),
  data = running, family = gaussian(),
  prior_intercept = normal(100, 10),
  prior = normal(2.5, 1),
  prior_aux = exponential(1, autoscale = TRUE),
  prior_covariance = decov(reg=1, conc=1,
                           shape=1, scale=1),
  chains = 2, iter = 1000, seed = 84735,
  adapt_delta = 0.99, refresh = 0)

bayesplot::mcmc_trace(as.array(m_slope), pars = "age")
summary(m_slope)$stan_summary |> head()
\end{lstlisting}
\textbf{Check:} $\widehat R$, effective sample size, divergences and trace plots; increase iterations/chains and tune if needed.
\end{frame}

\begin{frame}[fragile]{Check predictions, then compare models}
\begin{lstlisting}
library(bayesplot)
library(loo)
pp_check(m_int, nsamples = 40)
pp_check(m_slope, nsamples = 40)
loo_int <- loo(m_int)
loo_slope <- loo(m_slope)
loo_compare(loo_int, loo_slope)
\end{lstlisting}
Book's longer fits: ELPD difference $\approx1.6$ in favor of varying slopes, SE $\approx1.2$; modest evidence, not a decisive win.\par
\textbf{Task:} Give one reason to prefer simpler intercepts, and one reason to prefer richer slopes.
\end{frame}

\begin{frame}{MCQ 4 --- Who are we predicting for?}
You want to assess predictions for \emph{entirely new runners}. What should the test folds hold out?
\begin{enumerate}
\item[A.] Random individual races from all runners
\item[B.] Whole runners, with all of their races
\item[C.] Only races where net time is missing
\item[D.] Random posterior draws
\end{enumerate}
\par \vspace{4mm}
How would your split change for predicting the next race of a runner already observed?
\end{frame}

\begin{frame}{Build a prediction someone could use}
\stage{Scenario}{Predict net race time at age 61 for sampled runners 1 and 10 and for a new runner, ``Miles.''}
\begin{enumerate}
\item Identify known and unknown runner effects.
\item Specify whether you need an expected time or an individual future race time.
\item Predict which runner will have the widest interval and why.
\end{enumerate}
\textbf{Key distinction:} Existing runner predictions condition on information from that runner; new-runner predictions must also integrate between-runner heterogeneity.
\end{frame}

\visual{A new runner means more uncertainty}{ch17_predictions.png}{Illustrative intervals, not fitted results: a new runner has an unknown baseline.}

\begin{frame}[fragile]{Coding lab: predictions for new and existing runners}
\begin{lstlisting}
new_races <- data.frame(
  runner = factor(c("1", "10", "Miles"),
                  levels = c(levels(running$runner), "Miles")),
  age = 61)
# rstanarm: sample unseen runner effects from population
set.seed(84735)
yrep <- posterior_predict(m_int, newdata = new_races,
                          allow_new_levels = TRUE,
                          sample_new_levels = "gaussian")
apply(yrep, 2, median)
apply(yrep, 2, quantile, probs = c(.10, .90))
\end{lstlisting}
\textbf{Compare:} medians and 80\% predictive intervals. Confirm your installed \texttt{rstanarm} version supports the new-level arguments.
\end{frame}

\begin{frame}{MCQ 5 --- Predicting a typical time or a real race?}
Which quantity is appropriate for a \emph{new race outcome}?
\begin{enumerate}
\item[A.] Only $\beta_{0j}+\beta_1x$, with all uncertainty removed
\item[B.] A posterior predictive draw incorporating residual scatter
\item[C.] The prior mean of $\beta_1$ alone
\item[D.] The group intercept standard deviation alone
\end{enumerate}
\end{frame}

\begin{frame}{Choose a model and explain why}
\begin{enumerate}
\item Why reject complete pooling for these running data?
\item Does slope heterogeneity justify the more complex model?
\item What do predictive checks and held-out validation assess differently?
\item What assumptions remain untested (linearity, residual structure, longitudinal dependence)?
\end{enumerate}
\vspace{3mm}\stage{Book conclusion}{Prefer the varying-intercept model: it captures clustering, while the extra varying slopes add little predictive value in this example.}
\end{frame}

\begin{frame}{MCQ 6 --- Is the more complex model worth it?}
The richer model's ELPD improves only slightly, with uncertainty comparable to the difference. The simplest defensible conclusion is:
\begin{enumerate}
\item[A.] The richer model is guaranteed better for new runners.
\item[B.] The pooled model must be selected because it has fewer parameters.
\item[C.] Prefer the varying-intercept model if the extra complexity lacks a clear benefit.
\item[D.] Posterior predictive checks no longer matter.
\end{enumerate}
\end{frame}

\appendix

\begin{frame}{Answer key --- MCQs 1--3}
\textbf{1: B.} Races by a given runner can be correlated; pooled independence is questionable.\par\vspace{3mm}
\textbf{2: A.} \texttt{(1 | runner)} permits varying intercepts; slopes remain shared. As $\sigma_0\to0$, the random-intercept component collapses to the pooled mean structure.\par\vspace{3mm}
\textbf{3: B.} Negative intercept--slope correlation means lower initial times (faster runners) are associated with steeper age trends.
\end{frame}

\begin{frame}{Answer key --- MCQs 4--6}
\textbf{4: B.} Hold out entire runners to evaluate generalization to unseen groups; a different split can be appropriate for future races by observed runners.\par\vspace{3mm}
\textbf{5: B.} A posterior predictive draw includes both parameter uncertainty and outcome variation.\par\vspace{3mm}
\textbf{6: C.} Complexity must earn its keep through fit, prediction, or scientific usefulness.
\end{frame}

\end{document}